\section{TEKNOLOGI}

\subsection{Daftar Teknologi}

Berikut adalah daftar lengkap teknologi yang digunakan dalam pengembangan Savora:

\begin{table}[H]
\centering
\caption{Daftar Teknologi Savora}
\label{tab:daftar_teknologi}
\small
\begin{tabularx}{\textwidth}{|l|X|X|}
\hline
\textbf{Kategori} & \textbf{Teknologi} & \textbf{Versi} \\
\hline
Bahasa Pemrograman & JavaScript, TypeScript & ES2024, 5.x \\
\hline
Frontend Framework & React.js & 18.x \\
\hline
Backend Framework & Node.js + Express & 20 LTS, 4.x \\
\hline
Database & PostgreSQL, Redis & 16.x, 7.x \\
\hline
ORM & Prisma & 5.x \\
\hline
Authentication & JWT, bcrypt & - \\
\hline
Real-time & Socket.IO & 4.x \\
\hline
Machine Learning & Python, scikit-learn & 3.12, 1.4 \\
\hline
API Design & RESTful API & - \\
\hline
Version Control & Git, GitHub & - \\
\hline
CI/CD & GitHub Actions & - \\
\hline
Hosting & AWS / Vercel & - \\
\hline
\end{tabularx}
\end{table}

\subsection{Alasan Pemilihan Teknologi}

\subsubsection{React.js untuk Frontend}
React.js dipilih sebagai \textit{frontend framework} karena:
\begin{itemize}[leftmargin=*]
    \item \textbf{Komponen Reusable} — Arsitektur berbasis komponen memungkinkan pembuatan UI yang konsisten dan mudah dipelihara.
    \item \textbf{Performa Tinggi} — Virtual DOM memberikan performa yang optimal untuk aplikasi dengan banyak interaksi.
    \item \textbf{Ekosistem Kaya} — Tersedia banyak library pendukung seperti React Router, Redux, dan Material-UI.
    \item \textbf{Komunitas Besar} — Mudah menemukan dokumentasi, tutorial, dan solusi untuk berbagai masalah.
\end{itemize}

\subsubsection{Node.js + Express untuk Backend}
Node.js dan Express dipilih karena:
\begin{itemize}[leftmargin=*]
    \item \textbf{Asynchronous} — Sangat cocok untuk aplikasi yang membutuhkan pemrosesan banyak request secara simultan.
    \item \textbf{JavaScript Full-stack} — Penggunaan bahasa yang sama di frontend dan backend mempercepat pengembangan.
    \item \textbf{Ringan dan Cepat} — Cocok untuk microservice dan API yang responsif.
    \item \textbf{NPM Ecosystem} — Akses ke ribuan package yang sudah teruji.
\end{itemize}

\subsubsection{PostgreSQL + Redis}
\begin{itemize}[leftmargin=*]
    \item \textbf{PostgreSQL} — Database relasional yang robust, mendukung ACID, dan memiliki fitur JSON untuk data semi-terstruktur.
    \item \textbf{Redis} — In-memory database untuk caching, session management, dan real-time features.
\end{itemize}

\subsubsection{Python + scikit-learn untuk Machine Learning}
Python dipilih untuk komponen \textit{machine learning} karena:
\begin{itemize}[leftmargin=*]
    \item \textbf{Ekosistem ML Terkaya} — scikit-learn, pandas, dan numpy menyediakan tools yang komprehensif untuk algoritma scoring dan analitik.
    \item \textbf{Integrasi Mudah} — Dapat diintegrasikan dengan backend Node.js melalui API atau \textit{microservice}.
    \item \textbf{Prototyping Cepat} — Memungkinkan eksperimen dan iterasi model dengan cepat.
\end{itemize}

\subsection{Arsitektur Teknologi}

\begin{figure}[H]
\centering
% Placeholder untuk diagram arsitektur
\fbox{\parbox{0.8\textwidth}{\centering\vspace{3cm}Diagram Arsitektur Sistem\\(Akan diganti dengan diagram aktual)\vspace{3cm}}}
\caption{Arsitektur High-Level Savora}
\label{fig:arsitektur}
\end{figure}

Alur komunikasi sistem:
\begin{enumerate}[leftmargin=*]
    \item \textbf{Client Layer} — Browser pengguna mengakses React.js application.
    \item \textbf{API Gateway} — Request diteruskan ke Express.js backend.
    \item \textbf{Application Layer} — Backend memproses business logic, autentikasi, dan validasi data.
    \item \textbf{Data Layer} — PostgreSQL menyimpan data persisten, Redis menyimpan cache dan session.
    \item \textbf{ML Service} — Python microservice menangani Food Eligibility Score calculation dan analitik dashboard.
    \item \textbf{Real-time Layer} — Socket.IO menangani notifikasi dan update status secara real-time.
\end{enumerate}
